\section*{第 \ref{ch:b3:organometallic-bonding} 章——有机金属化学：成键与配体}

\begin{solution}{exo:b3:organometallic-bonding:1}
$\ce{Fe(CO)5}$: $8 + 10 = 18$. $\ce{Mn2(CO)10}$: $7 + 10 + 1 = 18$ per Mn. $\ce{CpMo(CO)3H}$: $6 + 5 + 6 + 1 =
18$. $\ce{[PtCl3(C2H4)]-}$: $10 + 3 + 2 + 1 = 16$. $\ce{Cr($\eta^6$-C6H6)(CO)3}$: $6 + 6 + 6 = 18$. $\ce{Cp2ZrCl2}$:
$4 + 10 + 2 = 16$.
\end{solution}

\begin{solution}{exo:b3:organometallic-bonding:2}
Fe(0) $d^8$; Mn(0) $d^7$; Mo(II) $d^4$; Pt(II) $d^8$; Cr(0) $d^6$; Zr(IV) $d^0$.
\end{solution}

\begin{solution}{exo:b3:organometallic-bonding:3}
$\ce{[Mn(CO)6]+} > \ce{Cr(CO)6} > \ce{[V(CO)6]-}$：金属越富电子，
向 $\pi^*$ 的反馈键合越强，C--O 键越弱。
\end{solution}

\begin{solution}{exo:b3:organometallic-bonding:4}
乙烷（$\ce{Mn(CO)5}$ 和 $\ce{CpFe(CO)2}$ 与 \ce{CH3} \omterm{def:b3:organometallic-bonding:isolobal}{等瓣的}）；又是乙烷；四面体烷
$\ce{(CH)4}$ 中三个 CH 被 $\ce{Co(CO)3}$ 取代。
\end{solution}

\begin{solution}{exo:b3:organometallic-bonding:5}
面式（$C_{3v}$）：两个（A$_1$ + E）。经式（$C_{2v}$）：三个（2A$_1$ + B$_1$）。
\end{solution}

\begin{solution}{exo:b3:organometallic-bonding:6}
体积较大的膦使金属周围拥挤，解离可消除张力，
因此解离更快，低配位中间体也更有利。
\end{solution}

\begin{solution}{exo:b3:organometallic-bonding:7}
该铬配合物是\omterm{def:b3:organometallic-bonding:carbene}{费歇尔卡宾}（取代基为 OMe，Cr(0)，CO 配体）：胺
进攻亲电的\omterm{def:b3:organometallic-bonding:carbene}{卡宾}碳并取代 OMe。钽的
亚烷基是\omterm{def:b3:organometallic-bonding:carbene}{施罗克卡宾}，其亲核碳进攻
酮的羰基，生成烯烃与 Ta=O 单元，类似维蒂希反应。
\end{solution}

\begin{solution}{exo:b3:organometallic-bonding:8}
交叉式：无 $\delta$ 重叠，两个 $\delta$ 电子为非键，键级为 3。$\ce{[Re2Cl8]^{4-}}$：
$\sigma^2\pi^4\delta^2\delta^{\ast2}$，键级为 3。
\end{solution}

\begin{solution}{exo:b3:organometallic-bonding:9}
M$\cdots$H 距离短（衍射，中子衍射最好）；C--H 伸缩
振动波数低；核磁共振谱信号位于高场，单键 C--H 耦合减小。
\end{solution}

\begin{solution}{exo:b3:organometallic-bonding:10}
16 电子：平面正方形 $d^8$（$\ce{[PtCl4]^{2-}}$，威尔金森催化剂），其第 18 电子轨道
能量过高；早期金属的\omterm{def:b3:organometallic-bonding:metallocene}{茂金属}，如 $\ce{Cp2ZrCl2}$，
空间过于拥挤，无法再接纳配体。17 电子：$\ce{V(CO)6}$，因空间位阻不能二聚；19 电子：二茂钴，
多出的电子位于反键轨道中，易失去。
\end{solution}

\begin{solution}{exo:b3:organometallic-bonding:11}
CO 是强 $\pi$ 受体：它降低 $t_{2g}$ 能级，使 $\Delta_{\mathrm o}$ 很大，远高于
成对能：$t_{2g}^6$，抗磁性。$d$--$d$ 跃迁能量为 $\Delta_{\mathrm o}$，落在紫外区，
故不吸收可见光。
\end{solution}

\begin{solution}{exo:b3:organometallic-bonding:12}
$\pi$ 电子给予与向 $\pi^*$ 反馈键合都会降低 C--C 键级。在蔡斯盐
（Pt(II)，反馈键合适中）中，键长略有增加；低价、
富电子的金属反馈键合很强，配合物接近
\omterm{def:b3:organometallic-bonding:dcd}{金属环丙烷}极限，C--C 键接近单键。
\end{solution}

\begin{solution}{pb:b3:organometallic-bonding:1}
\textbf{1.} $\ce{Fe^{2+}}$ ($d^6$) $+ 2 \times 6 = 18$.
\textbf{2.} $8 + 2 \times 5 = 18$.
\textbf{3.} $9 + 10 = 19$.
\textbf{4.} $10 + 10 = 20$.
\textbf{5.} $18 - 1 = 17$.
\textbf{6.} Fe(II) $d^6$; Co(II) $d^7$; Ni(II) $d^8$; Fe(III) $d^5$.
\textbf{7.} $a_{1g}$ ($d_{z^2}$), $e_{1g}$ ($d_{xz}$, $d_{yz}$), $e_{2g}$ ($d_{xy}$, $d_{x^2-y^2}$).
\textbf{8.} $e_{1g}$ 对与充满电子的环 $e_{1g}$ 组合强烈重叠：成键
轨道主要位于环上，反键 $e_{1g}^\ast$ 主要位于金属上，
能量较高。
\textbf{9.} $e_{2g} \approx a_{1g} < e_{1g}^\ast$.
\textbf{10.} $e_{2g}^4a_{1g}^2$：无未成对电子。
\textbf{11.} $e_{1g}^\ast$ 中一个电子：一个未成对电子，$1.73\,\mu_B$。
\textbf{12.} 两个电子分占两个 $e_{1g}^\ast$ 轨道，自旋平行：两个未成对电子，$\sqrt{8} =
2.83\,\mu_B$.
\textbf{13.} $e_{2g}^3a_{1g}^2$：一个未成对电子。$e_{2g}^3$ 空穴给出轨道简并的
（E）态，并带有\omterm{def:b3:complex-spectra-magnetism:moment}{轨道贡献}。
\textbf{14.} 第 19 个电子位于反键轨道中，易被移去，
生成 18 电子的二茂钴正离子。
\textbf{15.} 被移去的电子来自近似非键的轨道，
结构几乎不变（\omterm{def:b3:complex-mechanisms:reorganisation}{重组能}小，见\cref{ch:b3:complex-mechanisms}）。
\textbf{16.} 该电对在多数溶剂中快速且可逆，化合物
可溶且稳定，其电势对溶剂依赖很小，因为
电荷分散在较大的离子上。
\textbf{17.} 20 个电子，其中两个为反键电子，它通过失去这些电子进行反应：
通过环滑移或脱去一个环来结合配体，或被氧化。
\textbf{18.} $\ce{CpFe(CO)2}$：$8 + 5 + 4 = 17$ 个电子，比 18 少一个，有一个前线轨道，与
\ce{CH3} 等瓣。二聚体 $\ce{[CpFe(CO)2]2}$ 即“乙烷”，含 Fe--Fe 键（实际上有两个 CO
配体桥连其间）。
\textbf{19.} $\ce{CpFe(CO)2-Mn(CO)5}$，含 Fe--Mn 键。
\textbf{20.} $\ce{PPh3}$ 比 CO 给电子能力强、$\pi$ 受体能力弱：金属向
其余 CO 反馈的电子更多，其伸缩频率降低。
\textbf{21.} 两个（A$_1$ + E）。
\textbf{22.} 环滑移至 $\eta^3$ 可释放相当于两个电子的配位能力，进入的
配体可以缔合方式结合，而金属电子数不超过 18。
\textbf{23.} \textbf{$\mu_{\text{唯自旋}}(\text{二茂镍}) = \sqrt{2 \times 4} = 2.83\,\mu_B$。}
\end{solution}
